\section{Introduzione}
Una domanda come «E in quel caso chi può presentare ricorso?» lascia indeterminati sia il procedimento sia la persona interessata. I turni precedenti consentono di ricostruire entrambi, ma possono anche introdurre termini riferiti a casi ormai esclusi. Un sistema di recupero deve scegliere quali elementi della storia preservare nella ricerca. Una volta fornito il passaggio pertinente, il generatore affronta un secondo problema: stabilire quali condizioni riportare nella risposta e quali affermazioni l'evidenza consenta effettivamente di sostenere.

RETECO (SemEval 2027, Task~1) studia il recupero condizionato dal tempo nel Track~1, costruito su TEMPO, e dal dialogo nel Track~2, costruito su RECOR \cite{reteco,recor}. Ci concentriamo sul secondo. Nel subtrack~2a ogni turno richiede una lista ordinata di documenti del dominio. Nel~2b i passaggi pertinenti sono forniti dagli organizzatori e il sistema produce una risposta in inglese. Questa scelta sperimentale permette di attribuire un cambiamento nella risposta al generatore anche quando la ricerca sul corpus è difficile. Il caso in cui la generazione dipende dai documenti recuperati appartiene al subtrack~2c ed è esterno al presente studio.

La domanda che poniamo è operativa: una seconda applicazione dello stesso strato del decoder modifica utilmente l'ordine dei candidati o il sostegno documentale delle risposte? Il costo aggiuntivo può essere giustificato per un compito e risultare superfluo per l'altro. Nel 2a misuriamo quindi sia la disponibilità di evidenze tra i candidati sia la loro posizione dopo il riordinamento. Nel 2b valutiamo separatamente il contributo dell'adattamento supervisionato e quello della profondità a inferenza. Il primo risultato empirico riguarda la scelta dei candidati; l'efficacia del loop costituisce l'ipotesi da verificare.
